% Complete Appendix A.4 reconstruction and all stated extensions.

\begin{annexitem}{annex:p15.signed-measure-examples}{thm:p15-signed-measure-examples}{Differences, density integrals and the sign of every subset}
\begin{proof}
Let \(\alpha=\nu_1-\nu_2\), where \(\nu_2(\Omega)<\infty\).
Every value is in \((-\infty,\infty]\), and the empty value is zero.
For pairwise disjoint \(E_i\), write \(a_i=\nu_1(E_i)\) and
\(b_i=\nu_2(E_i)\). Then \(\sum_i b_i\leq\nu_2(\Omega)<\infty\).
If \(\sum_i a_i<\infty\), every difference is finite and
\(\sum_i|a_i-b_i|\leq\sum_i(a_i+b_i)<\infty\).
Its sum is \(\nu_1(\bigcup_iE_i)-\nu_2(\bigcup_iE_i)\).
If \(\sum_i a_i=\infty\), then
\(\sum_i(a_i-b_i)^-\leq\sum_i b_i<\infty\), and the positive-part
sum must be infinite: otherwise the inequality
\(a_i\leq(a_i-b_i)^++b_i\) would make \(\sum_i a_i\) finite.
Thus the series has exactly the required meaning in this case as well.
This proves the first construction.

For the integral construction, the previously proved nonnegative
integration theorem makes
\(\nu_1(E)=\int_E f^+\,d\mu\) and
\(\nu_2(E)=\int_E f^-\,d\mu\) measures.
The latter is finite by hypothesis, and their difference is the defined
signed integral on each event. Apply the first construction.

If \(\mu(A\cap\{f<0\})=0\), every measurable subset of \(A\) has
zero negative-part integral and nonnegative signed integral, so \(A\)
is positive. Conversely take \(E=A\cap\{f<0\}\).
If \(A\) is positive, then
\[
                     0\leq\alpha(E)=-\int_E f^-\,d\mu\leq0.
\]
The zero-integral criterion makes \(f^-=0\) almost everywhere on \(E\).
Since it is strictly positive everywhere on \(E\), this implies
\(\mu(E)=0\).
The negative-set criterion follows in the same way with
\(E=A\cap\{f>0\}\). There the integral is nonnegative, and negativity
of \(A\) forces it to be zero, including exclusion of an infinite
positive integral. Its zero-integral criterion again gives \(\mu(E)=0\).
The reverse implication follows by integrating a nonpositive function
outside a null set. These arguments concern every measurable subset,
not merely the total integral on \(A\).
\end{proof}
\end{annexitem}

\begin{annexitem}{annex:p15.signed-set-closure}{thm:p15-signed-set-closure}{Countable unions without monotonicity of a signed measure}
\begin{proof}
Any measurable subset of a positive set has all its measurable subsets
inside that same positive set, proving the first assertion.
Given positive \(A_n\), disjointify them by
\(D_n=A_n\setminus\bigcup_{j<n}A_j\).
For measurable \(E\subseteq\bigcup_nA_n\), the sets \(E\cap D_n\)
are disjoint and positive. Countable additivity gives
\(\alpha(E)=\sum_n\alpha(E\cap D_n)\geq0\), proving positivity
of the union.

The subset argument for negative sets is identical.
Disjointifying their union gives only nonpositive terms in the
corresponding series. Its sum cannot be \(+\infty\), since that
would require an infinite positive-part sum. It cannot be \(-\infty\),
since no signed-measure value is allowed to be \(-\infty\).
It is therefore a finite nonpositive number, by the signed-additivity
axiom. This proves negativity of the union. Finally any negative set's
value belongs to \((-\infty,0]\), and so is finite. No monotonicity of
\(\alpha\) on arbitrary events was assumed.
\end{proof}
\end{annexitem}

\begin{annexitem}{annex:p15.negative-subset}{thm:p15-negative-subset}{Successive removal and the minimal reciprocal thresholds}
\begin{proof}
The number \(\alpha(E)<0\) is finite. Every measurable subset of \(E\)
also has finite signed measure: if one had value \(+\infty\), finite
additivity with its complement in \(E\), whose value is greater than
\(-\infty\), would give \(\alpha(E)=+\infty\).
Set \(F_1=E\). If it is negative, take \(F=E\).
Otherwise there is a measurable positive-valued subset of \(F_1\).
Choose the least integer \(n_1\geq1\) for which some
\(E_1\subseteq F_1\) has \(\alpha(E_1)\geq1/n_1\), and choose
one such \(E_1\).

At stage \(k\), put \(F_k=E\setminus\bigcup_{j<k}E_j\).
If \(F_k\) is negative, finite additivity gives
\[
                    \alpha(F_k)=\alpha(E)-\sum_{j<k}\alpha(E_j)
                                  \leq\alpha(E)<0,
\]
and it is a suitable output. Otherwise choose the least \(n_k\geq1\)
admitting \(E_k\subseteq F_k\) with \(\alpha(E_k)\geq1/n_k\).
All the removed sets are disjoint.

If this process never stops, let
\(F=E\setminus\bigcup_{k\geq1}E_k\).
The signed additivity axiom for this partition of the finite-valued
event \(E\) gives an absolutely convergent series:
\[
                         \alpha(E)=\alpha(F)+\sum_{k\geq1}\alpha(E_k).
\]
Thus the nonnegative sum on the right is finite and
\(\alpha(F)\leq\alpha(E)<0\).
If \(F\) were not negative, it would contain a measurable \(G\) with
\(\alpha(G)>0\). Choose \(N\geq1\) with \(\alpha(G)\geq1/N\).
Because \(G\subseteq F_k\) at every stage, the least-index choice
would give \(n_k\leq N\) for every \(k\). Hence every removed value
would be at least \(1/N\), contradicting finiteness of their sum.
Therefore \(F\) is negative. This proves both the terminating and
nonterminating cases of the source construction.
\end{proof}
\end{annexitem}

\begin{annexitem}{annex:p15.hahn-decomposition}{thm:p15-hahn-decomposition}{Attaining the infimum of negative-set values}
\begin{proof}
Let \(c\) be the infimum of the negative-set values. The empty set
is negative, so \(c\leq0\), while initially \(c=-\infty\) is possible.
If \(c>-\infty\), choose negative \(B_n\) with
\(\alpha(B_n)<c+1/n\). If \(c=-\infty\), choose them with
\(\alpha(B_n)\leq-n\). Set \(B=\bigcup_nB_n\).
It is negative by the closure lemma and so has finite value.
Since \(B\setminus B_n\) is negative, finite additivity gives
\[
                   \alpha(B)=\alpha(B_n)+\alpha(B\setminus B_n)
                                      \leq\alpha(B_n).
\]
The choice in the infinite-infimum case would make this finite number
at most \(-n\) for every \(n\), a contradiction.
Thus \(c\) is finite. The same inequality in the other case gives
\(\alpha(B)\leq c\), whereas the definition of the infimum gives
\(\alpha(B)\geq c\). Hence \(\alpha(B)=c\).

Let \(A=B^c\). If \(A\) were not positive, there would be measurable
\(E\subseteq A\) with \(\alpha(E)<0\).
The preceding lemma gives a negative \(F\subseteq E\) with
\(\alpha(F)<0\). The disjoint union \(B\cup F\) is negative and
has value \(c+\alpha(F)<c\), contradicting the infimum.
This proves the required partition.

For two Hahn partitions, \(A_1\cap B_2\) is simultaneously positive
and negative, since it is a subset of each. Every measurable subset
therefore has value zero. The same holds for \(A_2\cap B_1\).
These are exactly the two pieces on which the positive parts of the
partitions disagree, proving the stated uniqueness qualification.
\end{proof}
\end{annexitem}

\begin{annexitem}{annex:p15.jordan-decomposition}{thm:p15-jordan-decomposition}{The mutually singular measures and the partition formula}
\begin{proof}
Choose a Hahn partition \((A,B)\). Its positive and negative restrictions
are nonnegative countably additive measures, because signed additivity
on disjoint subsets of \(A\) or of \(B\) has terms of only one sign.
The negative restriction has total mass \(-\alpha(B)<\infty\).
They vanish on \(B\) and \(A\), respectively, so are mutually singular.
Finite additivity on \(E=(E\cap A)\mathbin{\dot\cup}(E\cap B)\)
gives \(\alpha(E)=\alpha^+(E)-\alpha^-(E)\).

Suppose also \(\alpha=\nu_1-\nu_2\), with \(\nu_1,\nu_2\)
mutually singular and the difference defined in the stated convention.
Choose \(D\) with \(\nu_1(D)=0,\nu_2(D^c)=0\).
Then \(\alpha(D)=-\nu_2(D)\) cannot be \(-\infty\), so \(\nu_2\)
is finite. The pair \((D^c,D)\) is a Hahn partition. On its two pieces,
\[
                   \nu_1(E)=\alpha(E\cap D^c),\qquad
                   \nu_2(E)=-\alpha(E\cap D).
\]
The disagreements with \((A,B)\) are signed-null by Hahn uniqueness.
Decomposing each restriction into its common part and the two
signed-null differences shows
\(\alpha(E\cap D^c)=\alpha(E\cap A)\) and
\(\alpha(E\cap D)=\alpha(E\cap B)\).
Thus \(\nu_1=\alpha^+\) and \(\nu_2=\alpha^-\).

For \(F\subseteq E\),
\(\alpha(F)=\alpha^+(F)-\alpha^-(F)\leq\alpha^+(E)\).
Choosing \(F=E\cap A\) attains \(\alpha^+(E)\), even when it is infinite.
Similarly \(\alpha(F)\geq-\alpha^-(E)\), with equality at \(E\cap B\).
This proves the supremum and infimum formulas.

For every event \(D\), the extended scalar inequality
\[
                         |\alpha(D)|\leq\alpha^+(D)+\alpha^-(D)
\]
holds; if \(\alpha^+(D)=\infty\), both the relevant sides are infinite.
Adding over a finite measurable partition of \(E\) gives the upper
bound \(\lvert\alpha\rvert(E)\).
The two-part partition \(E\cap A,E\cap B\), omitting or retaining
empty parts, attains that bound. This proves the variation formula.
An event is signed-null exactly when both extrema are zero, equivalently
when its variation is zero.
Finally \(\alpha^-(\Omega)\) is always finite, so
\(\alpha(\Omega)<\infty\) exactly when \(\alpha^+(\Omega)<\infty\).
In that case every event has finite value and total variation is finite.
Conversely a finite signed measure has finite value at \(\Omega\),
which proves the final equivalence.
\end{proof}
\end{annexitem}

\begin{annexitem}{annex:p15.signed-cdf-variation}{thm:p15-signed-cdf-variation}{Cumulative functions and the full reverse variation inequality}
\begin{proof}
For a finite measure, its masses on \((-\infty,x]\) are bounded and
nondecreasing. Continuity from above gives right continuity, and
continuity from above on a sequence of left half-lines tending to the
empty set gives limit zero at \(-\infty\).
Apply these facts to the two finite Jordan measures to obtain the
stated properties and \(H=F-G\).

Conversely the Stieltjes theorem constructs finite measures \(\mu_F,\mu_G\)
from the two specified functions. Because their limits at \(-\infty\)
are zero, their actual cumulative functions are exactly \(F,G\).
Their difference is a finite signed measure with cumulative function
\(F-G\). To prove uniqueness, suppose \(\alpha,\beta\) have that
same cumulative function. The finite positive measures
\(\alpha^++\beta^-\) and \(\beta^++\alpha^-\) then have identical
cumulative functions. Their total masses agree by letting \(x\to\infty\).
If that mass is zero both vanish; otherwise divide each by the common
positive total and apply the established uniqueness of probability laws
from their cumulative functions. The two measures agree, and subtracting
the finite terms gives \(\alpha=\beta\).

For an interval \(I=(a,b]\), each partition sum in the theorem is
\(\sum_j|\alpha((t_{j-1},t_j])|\), which is at most
\(\lvert\alpha\rvert(I)\) by the Jordan variation formula.
We prove the reverse inequality explicitly.
Let \(\tau\) be the finite measure \(\lvert\alpha\rvert\) restricted
to \(I\), and let \(\mathcal A_I\) be the algebra of finite unions of
half-open intervals contained in \(I\), with endpoints including \(a,b\)
when needed.
Every Borel subset \(E\) of \(I\) can be approximated by \(U\in\mathcal A_I\)
so that \(\tau(E\mathbin{\triangle}U)<\varepsilon\).
Here is the required approximation argument.
The class of sets with this property for every positive \(\varepsilon\)
contains the interval algebra and is closed under complements in \(I\),
because symmetric difference is preserved by complement.
It is closed under countable unions: first approximate a union by its
first finitely many members using continuity from below and finiteness
of \(\tau(I)\), then approximate each member with errors whose finite sum
is less than the remaining allowance. The symmetric difference of the
two finite unions is contained in the union of those errors.
This class is therefore a sigma-algebra containing the intervals,
which generate the Borel subsets of \(I\).

Take a Hahn positive set \(A\), and choose \(U\in\mathcal A_I\)
approximating \(A\cap I\) within \(\varepsilon\).
Refine all endpoints of \(U\), together with \(a,b\), into one finite
partition of \(I\). Each partition interval lies in \(U\) or its
complement in \(I\), so its variation sum is at least
\[
 \begin{aligned}
 \alpha(U)-\alpha(I\setminus U)
   &=\lvert\alpha\rvert(I)
           -2\bigl(\alpha^+(I\setminus U)+\alpha^-(U)\bigr)\\
   &\geq\lvert\alpha\rvert(I)-2\tau(U\mathbin{\triangle}(A\cap I))\\
   &>\lvert\alpha\rvert(I)-2\varepsilon .
 \end{aligned}
\]
The middle inequality uses that the Jordan measures are concentrated
on the respective Hahn sets: only misclassified points contribute.
Let \(\varepsilon\) decrease to zero and combine with the upper bound.
The increments use \((t_{j-1},t_j]\), which excludes any initial atom
at \(a\) exactly as asserted.

It remains to show that the Stieltjes measures of a representing pair
need not be the Jordan components. Take any function \(F\) with the
listed properties that is not identically zero, and put \(G=F\).
The signed measure \(\mu_F-\mu_G\) has cumulative function \(F-G=0\),
so by the uniqueness proved above it is the zero signed measure. For it,
\((\mathbb R,\varnothing)\) is a Hahn partition, so both Jordan
components vanish. Yet \(\mu_F((-\infty,x])=F(x)>0\) at any \(x\)
with \(F(x)\neq0\), since \(F\) is nonnegative, so \(\mu_F=\mu_G\)
is not the zero measure. The exact relation between a representing pair
and the Jordan components is \cref{thm:p15-w02b-p15-005-stieltjes-excess}.
\end{proof}
\end{annexitem}

\begin{annexitem}{annex:p15.lebesgue-decomposition}{thm:p15-lebesgue-decomposition}{Maximal dominated density and the sigma-finite assembly}
\begin{proof}
\paragraph{Finite measures and the maximal integral.}
First suppose \(\mu(\Omega),\nu(\Omega)<\infty\).
Let \(\mathcal G\) consist of all measurable \(g\geq0\) satisfying
\(\int_Eg\,d\mu\leq\nu(E)\) for every measurable \(E\).
It contains zero. If \(g,h\in\mathcal G\), split any event into
its intersections with \(\{g>h\}\) and \(\{g\leq h\}\).
Adding the two domination inequalities shows \(\max(g,h)\in\mathcal G\).
Put
\[
                  \kappa=\sup_{g\in\mathcal G}\int g\,d\mu
                                  \leq\nu(\Omega)<\infty.
\]
Choose \(g_n\in\mathcal G\) with integral greater than \(\kappa-1/n\);
when \(\kappa=0\), the zero function satisfies this requirement.
Let \(h_n=\max(g_1,\ldots,g_n)\) and \(h=\sup_nh_n\).
Monotone convergence on every event gives \(h\in\mathcal G\).
Since \(\int h_n\geq\int g_n>\kappa-1/n\), and each integral is
at most \(\kappa\), another monotone limit gives \(\int h\,d\mu=\kappa\).

Define \(\nu_{\mathrm r}(E)=\int_Eh\,d\mu\), a finite measure.
The residual \(\nu_{\mathrm s}(E)=\nu(E)-\nu_{\mathrm r}(E)\)
is nonnegative on every event. It is a measure: for a disjoint union,
the two finite nonnegative sums for \(\nu,\nu_{\mathrm r}\) may be
subtracted term by term, yielding the sum of the residuals.
Thus \(\nu=\nu_{\mathrm r}+\nu_{\mathrm s}\).

\paragraph{Singularity by the source Hahn construction.}
For each integer \(n\geq1\), take a Hahn partition
\((A_n,B_n)\) for the finite signed measure
\(\nu_{\mathrm s}-\mu/n\).
Its positivity on every subset of \(A_n\) gives, for any event \(E\),
\[
 \int_E(h+n^{-1}\mathbf1_{A_n})\,d\mu
   =\nu_{\mathrm r}(E)+n^{-1}\mu(E\cap A_n)
   \leq\nu_{\mathrm r}(E)+\nu_{\mathrm s}(E\cap A_n)
   \leq\nu(E).
\]
Hence this larger function belongs to \(\mathcal G\).
Its integral is \(\kappa+\mu(A_n)/n\), so maximality forces
\(\mu(A_n)=0\).
Let \(A=\bigcup_nA_n\). Then \(\mu(A)=0\), and \(A^c\subseteq B_n\)
for every \(n\). Negativity gives
\[
                      \nu_{\mathrm s}(A^c)\leq\mu(A^c)/n
                                             \leq\mu(\Omega)/n .
\]
Letting \(n\to\infty\) proves \(\nu_{\mathrm s}(A^c)=0\).
This uses only the reciprocal sequence of Hahn partitions actually
chosen; no unchosen intermediate value of the parameter is needed.
The integral part vanishes on every \(\mu\)-null set, so is absolutely
continuous with respect to \(\mu\).

\paragraph{A common partition for sigma-finite measures.}
Now let both measures be sigma-finite. Choose measurable covers
\(\Omega=\bigcup_iU_i=\bigcup_jV_j\) with \(\mu(U_i)<\infty\)
and \(\nu(V_j)<\infty\). Enumerate the countably many intersections
\(U_i\cap V_j\), and disjointify them to obtain a measurable partition
\((D_m)_{m\geq1}\) with both measures finite on each member.
Empty members cause no difficulty.
Apply the finite construction on each \(D_m\), with its trace sigma-algebra.
It gives a measurable density \(h_m\) there and a residual \(\sigma_m\)
concentrated on \(A_m\subseteq D_m\), where \(\mu(A_m)=0\).
Extend each density by zero outside \(D_m\) and set
\[
 g=\sum_mh_m\mathbf1_{D_m},\quad
 \nu_{\mathrm r}(E)=\int_Eg\,d\mu,\quad
 \nu_{\mathrm s}(E)=\sum_m\sigma_m(E\cap D_m).
\]
The piecewise function is measurable. Countable additivity and
nonnegative integral summation give
\(\nu_{\mathrm r}(E)+\nu_{\mathrm s}(E)=\sum_m\nu(E\cap D_m)=\nu(E)\).
This assembly never subtracts two infinite values.
The union \(A=\bigcup_mA_m\) has \(\mu(A)=0\) and
\(\nu_{\mathrm s}(A^c)=0\).
Also \(\int_{D_m}g\,d\mu\leq\nu(D_m)<\infty\), so \(g\) is finite
almost everywhere on every \(D_m\), and hence on \(\Omega\).
It may be set to zero on its measurable exceptional null set.

\paragraph{Uniqueness of both measures.}
Suppose \(\nu=\rho_1+\sigma_1=\rho_2+\sigma_2\), with
\(\rho_i\ll\mu\) and \(\sigma_i\perp\mu\).
Choose \(\mu\)-null supports \(A_i\) of the two singular measures
and let \(A=A_1\cup A_2\).
On \(A^c\), both singular measures vanish, so
\(\rho_1(E\cap A^c)=\nu(E\cap A^c)=\rho_2(E\cap A^c)\).
On \(A\), both regular measures vanish by absolute continuity.
Thus \(\rho_1=\rho_2\).
Each singular measure is concentrated on \(A\), where its value on
\(E\cap A\) is exactly \(\nu(E\cap A)\); hence \(\sigma_1=\sigma_2\)
as well. No cancellation of infinite totals occurs in this proof.
\end{proof}
\end{annexitem}

\begin{annexitem}{annex:p15.radon-nikodym}{thm:p15-radon-nikodym}{Eliminating the singular part and proving density uniqueness}
\begin{proof}
Use the proved decomposition \(\nu=\nu_{\mathrm r}+\nu_{\mathrm s}\),
and choose \(A\) with \(\mu(A)=0,\nu_{\mathrm s}(A^c)=0\).
Absolute continuity gives \(\nu(A)=0\), so
\(\nu_{\mathrm s}(A)\leq\nu(A)=0\). The singular measure is zero
on both \(A,A^c\), and hence identically zero.
The constructed regular density therefore represents all of \(\nu\).
Its almost-everywhere finiteness was proved in the sigma-finite assembly.

For uniqueness, let \(g,h\) be any nonnegative measurable densities.
Choose a countable cover \(D_m\) with \(\nu(D_m)<\infty\).
Then both integrals on each \(D_m\) are finite, so \(g,h\) are finite
\(\mu\)-a.e. there. Their union of exceptional sets is a measurable
\(\mu\)-null set; replace both functions by zero on that set.
For \(k,m\geq1\), put \(E_{km}=D_m\cap\{g\geq h+1/k\}\).
The two restricted integrals are finite and equal to \(\nu(E_{km})\).
Signed integral linearity therefore yields
\[
            0=\int_{E_{km}}(g-h)\,d\mu\geq\mu(E_{km})/k .
\]
Thus \(\mu(E_{km})=0\). The union over \(k,m\) is \(\{g>h\}\);
it is null. Interchanging \(g,h\) treats \(\{h>g\}\).
This proves uniqueness including the previously excluded infinite-value
sets. Measurable null-set changes preserve all integrals.
Finally an integral of a nonnegative measurable function over a
\(\mu\)-null set is zero, proving the converse absolute-continuity assertion.
\end{proof}
\end{annexitem}

\begin{annexitem}{annex:p15.density-calculus}{thm:p15-density-calculus}{Integral transport and the derivative identities}
\begin{proof}
\paragraph{Change of measure.}
For \(f=\mathbf1_E\), the equality is the defining density identity.
For a nonnegative simple \(f\), it follows by finite linearity.
For a general nonnegative measurable \(f\), take increasing nonnegative
simple approximations \(f_n\uparrow f\). Then \(f_ng\uparrow fg\),
with the prescribed convention at zero and infinity, and monotone
convergence on both sides proves
\(\int f\,d\nu=\int fg\,d\mu\).
For signed \(f\), apply this separately to \(f^+,f^-\).
Define the signed product to be zero where \(g=0\), and otherwise
use its ordinary signed extended value. Its positive and negative
parts are respectively \(f^+g,f^-g\); they cannot be simultaneously
positive at a point. Thus the two signed integrals have exactly the
same part integrals, and are defined together whenever at least one
of those integrals is finite.

\paragraph{Sums and scalar multiples.}
Write \(g_i=d\nu_i/d\mu\). On every event,
\[
                 (\nu_1+\nu_2)(E)=\int_E(g_1+g_2)\,d\mu.
\]
The sum measure is sigma-finite: intersect a countable finite-measure
cover for \(\nu_1\) with one for \(\nu_2\).
It is absolutely continuous, so uniqueness identifies the displayed
density with its derivative. A nonnegative finite scalar \(c\)
gives \((c\nu)(E)=\int_Ecg\,d\mu\). At \(c=0\), both measures and
densities are zero; at positive \(c\), sigma-finiteness is preserved.
This proves the scalar rule as well.

\paragraph{Chain rule and versions on null sets.}
Let \(g=d\nu/d\mu\) and \(h=d\pi/d\nu\).
Absolute continuity is transitive, because
\(\mu(E)=0\) implies \(\nu(E)=0\), which implies \(\pi(E)=0\).
The change-of-measure identity, applied to \(h\mathbf1_E\), gives
\[
                         \pi(E)=\int_Eh\,d\nu=\int_Ehg\,d\mu.
\]
Thus uniqueness identifies \(hg\) with \(d\pi/d\mu\).
If a version of \(h\) is changed on a measurable \(\nu\)-null set \(N\),
then \(\int_Ng\,d\mu=\nu(N)=0\).
Hence \(g=0\) \(\mu\)-a.e. on \(N\), and the product remains
unchanged \(\mu\)-a.e., using zero times infinity equals zero.
Changes in \(g\) on a \(\mu\)-null set likewise have no effect.
This explains the null-set meaning of the chain identity rather than
assuming the two derivatives have identical exceptional sets.

\paragraph{Reciprocals.}
Now suppose also \(\mu\ll\nu\).
The set \(Z=\{g=0\}\) has \(\nu(Z)=\int_Zg\,d\mu=0\), and thus
\(\mu(Z)=0\). The set where \(g=\infty\) is \(\mu\)-null by the
sigma-finite representation theorem. Both are also \(\nu\)-null.
Define \(h=1/g\) on \(\{0<g<\infty\}\) and zero elsewhere.
It is measurable and satisfies \(hg=1\) \(\mu\)-a.e.
Therefore
\[
                       \int_Eh\,d\nu=\int_Ehg\,d\mu=\mu(E)
\]
for every event, including events of infinite mass.
Uniqueness for the reverse derivative gives \(h=d\mu/d\nu\)
\(\nu\)-a.e.

Finally choose \(N\) with \(\mu_2(N)=0,\nu(N^c)=0\).
If \(\mu_1\ll\mu_2\), then \(\mu_1(N)=0\), so the same
set proves \(\mu_1\perp\nu\). This last argument uses neither
finiteness nor a density theorem.
\end{proof}
\end{annexitem}

\begin{annexitem}{annex:p15.finite-signed-density}{thm:p15-finite-signed-density}{Applying representation to the two finite Jordan parts}
\begin{proof}
If \(\mu(E)=0\), every measurable \(F\subseteq E\) is also
\(\mu\)-null and has \(\alpha(F)=0\).
The Jordan supremum and infimum formulas then give
\(\alpha^+(E)=\alpha^-(E)=0\).
Both finite Jordan measures are therefore absolutely continuous
with respect to \(\mu\).
The Radon--Nikodym theorem gives densities \(g_+,g_-\geq0\).
Their integrals over the whole space are the finite Jordan masses,
so they are integrable and finite almost everywhere.
Choose real measurable versions by setting them to zero on the
union of their measurable exceptional null sets.

For a Hahn partition \((A,B)\), the positive Jordan measure vanishes
on \(B\), and the negative one vanishes on \(A\).
Consequently \(g_+=0\) almost everywhere on \(B\) and \(g_-=0\)
almost everywhere on \(A\).
Set \(f=g_+-g_-\). It is real, measurable and integrable; the
disjoint-support facts give \(f^+=g_+\) and \(f^-=g_-\) almost
everywhere. Signed linearity yields
\[
                    \int_Ef\,d\mu
                       =\alpha^+(E)-\alpha^-(E)=\alpha(E).
\]
Adding the part densities gives the variation density \(|f|\)
and its total-mass formula.

If \(h\) is another integrable real density for \(\alpha\), then
on \(E_k=\{f\geq h+1/k\}\), both restricted integrals are finite
and equal. Thus
\[
                     0=\int_{E_k}(f-h)\,d\mu\geq\mu(E_k)/k.
\]
Every \(E_k\) is null, and their union is \(\{f>h\}\).
Interchanging the functions gives equality almost everywhere.
This proves all signed representation and uniqueness assertions.
\end{proof}
\end{annexitem}

\begin{annexitem}{annex:p15.finite-uniform-absolute-continuity}{thm:p15-finite-uniform-absolute-continuity}{A summable-error contradiction}
\begin{proof}
The epsilon-delta property immediately implies \(\nu\ll\mu\):
if \(\mu(E)=0\), apply it for every \(\varepsilon>0\) to obtain
\(\nu(E)=0\).

Conversely suppose \(\nu\ll\mu\), \(\nu(\Omega)<\infty\), but
the property fails. Some \(\varepsilon_0>0\) then admits, for every
positive \(\delta\), an event with \(\mu\)-mass less than \(\delta\)
and \(\nu\)-mass at least \(\varepsilon_0\).
Choose \(E_n\) with
\(\mu(E_n)<2^{-n}\) and \(\nu(E_n)\geq\varepsilon_0\).
Let \(F_N=\bigcup_{n\geq N}E_n\). Countable subadditivity gives
\(\mu(F_N)\leq\sum_{n\geq N}2^{-n}\to0\).
For \(F=\bigcap_NF_N\), monotonicity implies \(\mu(F)=0\),
so \(\nu(F)=0\).
The \(F_N\) decrease, and their \(\nu\)-masses are finite because
\(\nu\) is finite. Continuity from above therefore gives
\(\nu(F_N)\downarrow0\).
But \(E_N\subseteq F_N\) gives \(\nu(F_N)\geq\varepsilon_0\)
for every \(N\), a contradiction.
\end{proof}
\end{annexitem}

\begin{annexitem}{annex:p15.cantor-singular-law}{thm:p15-cantor-singular-law}{The fair-digit measure and every middle-third plateau}
\begin{proof}
For a word \(b\) of length \(n\), its two successor intervals are
\[
 [c_b,c_b+3^{-(n+1)}],\qquad
 [c_b+2\cdot3^{-(n+1)},c_b+3^{-n}].
\]
Thus they retain the two outer thirds of \(I_b\) and are disjoint.
Induction proves that the \(2^n\) intervals at stage \(n\) are
pairwise disjoint, have length \(3^{-n}\), and that \(C_{n+1}\subseteq C_n\).
They are closed Borel sets. Lebesgue lengths therefore give
\(\lambda(C_n)=(2/3)^n\); since \(C\subseteq C_n\) for every \(n\),
\(\lambda(C)=0\).

The partial sums of \(X\) are measurable and increase to a limit
bounded by \(\sum_{j\geq1}2/3^j=1\).
Hence \(X\) is a measurable real random variable.
If its first \(n\) bits equal \(b\), its remaining nonnegative tail
is at most \(3^{-n}\), so \(X\in I_b\).
Every outcome has some length-\(n\) prefix, and these stage intervals
are pairwise disjoint. Thus, conversely, membership in \(I_b\)
forces that prefix. Independence and fairness give
\[
                    \eta(I_b)
                        =\mathbb P(B_1=b_1,\ldots,B_n=b_n)=2^{-n}.
\]
Also every outcome belongs to every \(C_n\), so \(\eta(C)=1\).
If \(x\in C\), it lies in one stage-\(n\) interval for every \(n\),
and hence \(\eta(\{x\})\leq2^{-n}\) for all \(n\).
If \(x\notin C\), it has zero mass because \(\eta(C)=1\).
Therefore every singleton is null.
Any countable union of them is null as well, proving that the law
is not concentrated on a countable set.

The cumulative-distribution theorem now shows that the cumulative
function is continuous: its possible jumps are precisely its
singleton masses, all of which are zero.
It vanishes for \(x<0\) and equals one for \(x\geq1\) by the range
of \(X\); the null mass at zero gives the claimed value there too.

For the complete plateau formula, fix a binary word
\(b=(b_1,\ldots,b_{n-1})\), including the empty word at \(n=1\).
The middle third removed from its interval is
\[
               (c_b+3^{-n},\ c_b+2\cdot3^{-n}).
\]
The stage prefixes less than \(b\) in binary order contribute
probability \(\sum_{j=1}^{n-1}b_j2^{-j}\).
To verify this count, at the first position \(j\) where \(b_j=1\),
choosing the smaller bit 0 and retaining the earlier prefix has
probability \(2^{-j}\), and these first-difference events are
disjoint and exhaust the smaller prefixes.
Within prefix \(b\), the lower successor interval contributes
\(2^{-n}\); no point of the open removed interval has any mass.
The two endpoints are singletons of zero mass. Therefore throughout
the closed plateau interval the cumulative value is
\[
                    \sum_{j=1}^{n-1}b_j2^{-j}+2^{-n}.
\]
The empty prefix gives \(1/2\); the two one-bit prefixes give
\(1/4\) and \(3/4\), exactly the stated first plateaus.

The removed open intervals are dense in \([0,1]\).
Indeed, an open interval of positive length cannot lie in every
\(C_n\), because the longest connected interval in \(C_n\) has
length \(3^{-n}\), and the components are separated.
Thus the continuous cumulative function is uniquely determined by
these plateau values together with its endpoint values.
It is the source's middle-thirds distribution.

Finally, \(\lambda(C)=0\) and \(\eta(C^c)=0\) prove
\(\eta\perp\lambda\).
If a nonnegative Lebesgue density \(f\) represented it, then
\(\eta(C)=\int_C f\,d\lambda=0\), contradicting \(\eta(C)=1\).
This proves the absence of a density despite continuity of the
cumulative function.
\end{proof}
\end{annexitem}

\begin{annexitem}{annex:p15.radon-nikodym-examples}{thm:p15-radon-nikodym-examples}{The base measure, cancellation and sigma-finiteness}
\begin{proof}
On \([0,1]\), the measure with density \(2x\) relative to Lebesgue
measure is absolutely continuous. The point mass at zero is singular,
because Lebesgue measure gives \(\{0\}\) mass zero.
Their sum is the displayed Lebesgue decomposition, unique by the theorem.
For \(\mu=\lambda+\delta_0\), define \(g(0)=1\) and \(g(x)=2x\)
for \(x>0\). Integrating against the two summands of \(\mu\) gives
\[
                     \int_Eg\,d\mu=\int_E2x\,d\lambda+\delta_0(E)
                                      =\nu(E).
\]
The equality follows first for simple functions and then by monotone
convergence, so does not presume a change-of-measure density for the sum.
In particular \(\nu\ll\mu\) and \(g\) is its derivative.
Changing \(g(0)\) would change the represented point mass, because
zero is no longer null for this base measure.

For \(f(x)=2x-1\), direct one-variable integration gives
\[
 \int_0^1f(x)\,dx=0,\quad
 \int_{1/2}^1(2x-1)\,dx=\frac14,\quad
 \int_0^{1/2}(1-2x)\,dx=\frac14 .
\]
The sign criterion supplies a Hahn partition at \(1/2\);
its allocation of that zero-density singleton is immaterial, as
\cref{thm:p15-w02b-p15-005-density-hahn} shows.
The two positive-part integrals are the Jordan masses, and their sum
is total variation \(1/2\). Thus the whole space has zero signed value
but is not signed-null.

For the counting-measure example, the only null Borel set of counting
measure is the empty set. Hence Lebesgue measure is absolutely
continuous with respect to it.
A set of finite counting measure is finite, so a countable union
of finite-measure sets is countable and cannot cover the uncountable
real line. Counting measure is not sigma-finite.
If a density \(g\geq0\) represented Lebesgue measure, its integral
on \(\{x\}\) would be \(g(x)\), since that singleton has counting
mass one. Its required value is zero for every real \(x\).
Thus \(g\) would vanish everywhere, giving zero mass on \([0,1]\)
instead of one. No such density exists.

On \((0,1)\), let \(\nu(E)=\int_E x^{-1}\,d\lambda\).
This is a measure and is absolutely continuous by the zero-integral
property of null sets. The sets \([1/n,1)\), \(n\geq2\), cover the
space and have finite mass \(\log n\), so \(\nu\) is sigma-finite.
For any \(a\in(0,1)\), monotone integration on \([\varepsilon,a)\)
as \(\varepsilon\downarrow0\) gives
\(\nu((0,a))=\infty\).
Given any proposed \(\delta>0\), choose \(0<a<\min(\delta,1)\).
Then \(\lambda((0,a))=a<\delta\), while \(\nu((0,a))=\infty\),
contradicting the uniform epsilon-delta property for, say,
\(\varepsilon=1\). This verifies the necessity of the finite-target
qualification in that particular conclusion.
\end{proof}
\end{annexitem}
